\begin{Solbook}{1}
 \par \noindent  \begin {enumerate} \item Điều kiện xác định: $x\ne 1$, $x\ne -1$. \begin {align*} \dfrac {x+2}{x-1}+\dfrac {x-3}{x+1}&=\dfrac {4x+2}{x^2-1}\\ (x+2)(x+1)+(x-3)(x-1)&=4x+2\\ 2x^2-5x+3&=0\\ (2x-3)(x-1)&=0. \end {align*} Suy ra $x=\dfrac {3}{2}$ hoặc $x=1$. \\ Vì $x=1$ không thỏa mãn điều kiện xác định nên phương trình có nghiệm $x=\dfrac {3}{2}$. \par \item Nhân cả hai vế với $6>0$, ta được \begin {align*} 2(2x-3)-3(x+4)&\le 5(x-2)\\ x-18&\le 5x-10\\ -8&\le 4x\\ x&\ge -2. \end {align*} Vậy bất phương trình có nghiệm $x\ge -2$. \par \item Dùng phương pháp thế, ta có \begin {align*} &\begin {cases} (x+2)(y-1)=xy-6,\\ (x-1)(y+3)=xy+4 \end {cases}\\[0.7em] \Leftrightarrow \quad &\begin {cases} x-2y=4,\\ 3x-y=7 \end {cases}\\[0.7em] \Leftrightarrow \quad &\begin {cases} x=4+2y,\\ 3x-y=7 \end {cases}\\[0.7em] \Leftrightarrow \quad &\begin {cases} x=4+2y,\\ 3(4+2y)-y=7 \end {cases}\\[0.7em] \Leftrightarrow \quad &\begin {cases} x=4+2y,\\ 5y=-5 \end {cases}\\[0.7em] \Leftrightarrow \quad &\begin {cases} x=2,\\ y=-1. \end {cases} \end {align*} Vậy hệ phương trình có nghiệm $(x;y)=(2;-1)$. \end {enumerate}  \hfill \qedEX 
\end{Solbook}
\begin{Solbook}{2}
 \par \noindent  \begin {enumerate} \setlength {\itemsep }{1em} \par \item Ta có \begin {align*} x+6\sqrt {x}+9&=\left (\sqrt {x}+3\right )^2,\\ x-4\sqrt {x}+4&=\left (\sqrt {x}-2\right )^2,\\ x-10\sqrt {x}+25&=\left (\sqrt {x}-5\right )^2. \end {align*} Vì $4\le x\le 25$ nên $2\le \sqrt {x}\le 5$. Do đó \begin {align*} A &=\left |\sqrt {x}+3\right | +\left |\sqrt {x}-2\right | -\left |\sqrt {x}-5\right |\\ &=\sqrt {x}+3+\sqrt {x}-2-\left (5-\sqrt {x}\right )\\ &=3\sqrt {x}-4. \end {align*} Vậy $A=3\sqrt {x}-4$. \par \item Với $x\ge 0$ và $x\ne 4$, ta có \begin {align*} B &=\left (\dfrac {\sqrt {x}+2}{\sqrt {x}-2} -\dfrac {\sqrt {x}-2}{\sqrt {x}+2}\right ) \cdot \dfrac {x-4}{8}\\ &=\dfrac {\left (\sqrt {x}+2\right )^2 -\left (\sqrt {x}-2\right )^2}{x-4} \cdot \dfrac {x-4}{8}\\ &=\dfrac {8\sqrt {x}}{x-4}\cdot \dfrac {x-4}{8}\\ &=\sqrt {x}. \end {align*} Vậy $B=\sqrt {x}$. \end {enumerate}  \hfill \qedEX 
\end{Solbook}
\begin{Solbook}{3}
 \par \noindent  Gọi $x$ (nghìn đồng) là giá một vé người lớn buổi sáng và $y$ (nghìn đồng) là giá một vé trẻ em buổi sáng, với $x>0$, $y>0$. \par Buổi sáng cho ta \[ 3x+2y=400. \tag {1} \] Buổi chiều, vé người lớn có giá $\dfrac {3}{4}x$ và vé trẻ em có giá $\dfrac {9}{10}y$, nên \[ 4\cdot \dfrac {3}{4}x+2\cdot \dfrac {9}{10}y=390, \] tức là \[ 3x+\dfrac {9}{5}y=390. \] Nhân hai vế với $5$: \[ 15x+9y=1950. \] Chia hai vế cho $3$: \[ 5x+3y=650. \tag {2} \] Nhân $(2)$ với $2$ và nhân $(1)$ với $3$: \begin {align*} 10x+6y&=1300, \tag {3}\\ 9x+6y&=1200. \tag {4} \end {align*} Lấy $(3)$ trừ $(4)$: \[ x=100. \] Thay vào $(1)$: \[ 3\cdot 100+2y=400 \Rightarrow 2y=100 \Rightarrow y=50. \] Cặp $(x;y)=(100;50)$ thỏa $x>0$, $y>0$. \par \emph {Kiểm tra.} Buổi chiều, vé người lớn giá $75$ nghìn đồng và vé trẻ em giá $45$ nghìn đồng. \[ 4\cdot 75+2\cdot 45=300+90=390, \] đúng với giả thiết. \par Vậy buổi sáng, một vé người lớn giá $100\,000$ đồng và một vé trẻ em giá $50\,000$ đồng.  \hfill \qedEX 
\end{Solbook}
\begin{Solbook}{4}
 \par \noindent  \begin {enumerate} \item Diện tích hình chữ nhật là \[ S_{\text {cn}}=(x+5)(x+1)=x^2+6x+5\ \text {(cm}^2\text {)}. \] Diện tích hình tròn bán kính $2$ cm là \[ S_{\text {tròn}}=\pi \cdot 2^2=4\pi \ \text {(cm}^2\text {)}. \] Diện tích phần tô đậm là \[ S=S_{\text {cn}}-S_{\text {tròn}}=x^2+6x+5-4\pi \ \text {(cm}^2\text {)}. \] \par \item Ta có \begin {align*} x^2+6x+5-4\pi &=32-4\pi \\ x^2+6x+5&=32\\ x^2+6x-27&=0\\ (x+9)(x-3)&=0. \end {align*} Suy ra $x=-9$ hoặc $x=3$. \par Vì $x>0$ nên $x=-9$ bị loại. Với $x=3$, hình chữ nhật có chiều dài $8$ cm và chiều rộng $4$ cm. Đường kính hình tròn bằng $4$ cm nên hình tròn tiếp xúc hai cạnh chiều rộng và nằm trọn trong hình chữ nhật. \par Vậy $x=3$ cm. \end {enumerate}  \hfill \qedEX 
\end{Solbook}
\begin{Solbook}{5}
 \par \noindent  \begin {enumerate} \item Gọi $R$ (m) là bán kính mép ngoài, $R>5$. Diện tích hình vành khuyên là \begin {align*} \pi R^2-\pi \cdot 5^2&=75\pi \\ R^2-25&=75\\ R^2&=100. \end {align*} Vì $R>0$ nên $R=10$. \par Vậy bán kính mép ngoài là $10$ m. \par \item Cung nhỏ $AB$ thuộc đường tròn bán kính $10$ m và chắn góc ở tâm $60^\circ $, nên \[ \ell _{\overset {\frown }{AB}} =\dfrac {\pi Rn}{180} =\dfrac {\pi \cdot 10\cdot 60}{180} =\dfrac {10\pi }{3}\ \text {(m)}. \] Phần đường chạy trong góc $\widehat {AOB}$ là hình quạt vành khuyên, bằng hiệu hai hình quạt cùng góc $60^\circ $: \begin {align*} S &=\dfrac {60}{360}\pi \cdot 10^2-\dfrac {60}{360}\pi \cdot 5^2\\ &=\dfrac {1}{6}\pi (100-25)\\ &=\dfrac {75\pi }{6} =\dfrac {25\pi }{2}\ \text {(m}^2\text {)}. \end {align*} Vậy cung nhỏ $AB$ dài $\dfrac {10\pi }{3}$ m và phần đường chạy trong góc $\widehat {AOB}$ có diện tích $\dfrac {25\pi }{2}$ m$^2$. \end {enumerate}  \hfill \qedEX 
\end{Solbook}
\begin{Solbook}{6}
 \par \noindent  \begin {center} \begin {tikzpicture}[scale=1.1, font=\footnotesize , line join=round, line cap=round] \coordinate (O) at (0,0); \coordinate (A) at (0,4); \coordinate (B) at (-1.732,1); \coordinate (C) at (1.732,1); \coordinate (H) at (0,1); \coordinate (M) at (-0.866,2.5); \coordinate (N) at (0.866,2.5); \coordinate (D) at (-0.684,1.879); \coordinate (I) at (1.020,2.5); \coordinate (T) at (0.560,1.372); \par \draw [thick] (O) circle (2); \par \draw [thick] (A)--(B) (A)--(C); \draw [thick] (B)--(C); \draw [thick] (O)--(A); \draw [thick] (M)--(I); \draw [thick] (D)--(I); \par \draw [dashed] (O)--(B) (O)--(C) (O)--(D) (O)--(I) (D)--(T); \par \draw ($(T)+(0.08,0.03)$)--++(-0.09,0.12)--++(-0.12,-0.09); \draw ($(H)+(-0.08,0)$)--++(0,0.08)--++(0.08,0); \par \fill (A) circle (1.2pt) node[above] {$A$}; \fill (O) circle (1.2pt) node[left] {$O$}; \fill (B) circle (1.2pt) node[left] {$B$}; \fill (C) circle (1.2pt) node[right] {$C$}; \fill (H) circle (1.2pt) node[right] {$H$}; \fill (M) circle (1.2pt) node[left] {$M$}; \fill (N) circle (1.2pt) node[above] {$N$}; \fill (D) circle (1.2pt) node[above left] {$D$}; \fill (I) circle (1.2pt) node[right] {$I$}; \fill (T) circle (1.2pt) node[right] {$T$}; \end {tikzpicture} \end {center} \par \begin {enumerate} \item Vì $AB$, $AC$ là hai tiếp tuyến kẻ từ $A$ đến đường tròn $(O)$ nên \[ AB=AC. \] Mặt khác, \[ OB=OC \] (vì cùng bằng bán kính của đường tròn $(O)$). \par Vậy $A$ và $O$ đều cách đều hai điểm $B$, $C$, suy ra đường thẳng $AO$ là đường trung trực của đoạn thẳng $BC$. \par Do đó \[ OA\perp BC. \] \par \item Từ câu a), ta có $OA\perp BC$. Lại do $AO$ là đường trung trực của $BC$ nên \[ HB=HC=\dfrac {BC}{2}. \tag {1} \] \par Xét hai tam giác vuông $\triangle AHB$ và $\triangle AOB$, ta có: \begin {itemize} \item $\widehat {AHB}=90^\circ $; \item $\widehat {ABO}=90^\circ $ (bán kính vuông góc với tiếp tuyến); \item $\widehat {BAH}=\widehat {BAO}$. \end {itemize} Suy ra \[ \triangle AHB \sim \triangle AOB. \] \par Do đó \[ \dfrac {AH}{AB}=\dfrac {AB}{AO} \Rightarrow AH\cdot AO=AB^2. \tag {2} \] \par Mà trong tam giác vuông $AOB$, ta có \[ AB^2=AO^2-OB^2=AO^2-R^2. \] Thay vào (2), được \[ AH\cdot AO=AO^2-R^2. \] Suy ra \[ AO(AO-AH)=R^2. \] Vì $O,H,A$ thẳng hàng và $H$ nằm giữa $O,A$ nên \[ AO-AH=OH. \] Do đó \[ OA\cdot OH=R^2. \tag {3} \] \par Xét tam giác vuông $OHB$, ta có \[ HB^2=OB^2-OH^2=R^2-OH^2. \] Theo (3), $R^2=OA\cdot OH$, nên \[ HB^2=OA\cdot OH-OH^2=OH(OA-OH). \] Mà $OA-OH=HA$, suy ra \[ HB^2=OH\cdot HA. \] Nhân hai vế với $4$, ta được \[ 4HB^2=4OH\cdot HA. \] Theo (1), \[ BC=2HB \Rightarrow BC^2=4HB^2. \] Vậy \[ 4OH\cdot HA=BC^2. \] \par \item Gọi $K=OA\cap MN$. \par Vì $M$, $N$ lần lượt là trung điểm của $AB$, $AC$ nên trong tam giác $ABC$ ta có \[ MN\parallel BC. \] Suy ra \[ OA\perp MN \tag {4} \] (do $OA\perp BC$ theo câu a)). \par Lại vì $MN\parallel BC$, nên $\triangle AMN \sim \triangle ABC$ với tỉ số đồng dạng \[ \dfrac {AM}{AB}=\dfrac 12. \] Do đó \[ \dfrac {AK}{AH}=\dfrac 12 \Rightarrow AK=\dfrac {AH}{2}. \] Suy ra $K$ là trung điểm của $AH$, nên \[ AK=KH=\dfrac {AH}{2}. \tag {5} \] \par \medskip \textbf {Chứng minh $OH\cdot OA=OT\cdot OI$.} \par Vì $ID$ là tiếp tuyến của $(O)$ tại $D$ nên \[ OD\perp ID. \] Suy ra tam giác $ODI$ vuông tại $D$. \par Lại có $T$ là hình chiếu của $D$ trên $OI$, nên \[ DT\perp OI. \] Vậy $DT$ là đường cao ứng với cạnh huyền $OI$ của tam giác vuông $ODI$. \par Theo hệ thức lượng trong tam giác vuông, ta có \[ OD^2=OT\cdot OI. \] Mà $OD=R$, nên \[ R^2=OT\cdot OI. \tag {6} \] \par Mặt khác, theo (3), \[ OH\cdot OA=R^2. \tag {7} \] \par Từ (6) và (7), suy ra \[ OH\cdot OA=OT\cdot OI. \] \par \medskip \textbf {Chứng minh $IA=ID$.} \par Do $I\in MN$ và $K\in MN$, từ (4) suy ra hai tam giác $AKI$ và $OKI$ đều vuông tại $K$. \par Áp dụng định lý Py-ta-go, ta có \[ IA^2=IK^2+AK^2, \tag {8} \] \[ OI^2=IK^2+OK^2. \tag {9} \] \par Lấy (9) trừ (8), được \[ OI^2-IA^2=OK^2-AK^2=(OK-AK)(OK+AK). \tag {10} \] \par Từ (5), vì $K$ là trung điểm của $AH$ và các điểm $O,H,K,A$ thẳng hàng theo thứ tự đó, nên \[ OK-AK=OH, \qquad OK+AK=OA. \] Thay vào (10), ta được \[ OI^2-IA^2=OH\cdot OA=R^2. \tag {11} \] \par Mặt khác, trong tam giác vuông $ODI$, ta có \[ ID^2=OI^2-OD^2=OI^2-R^2. \tag {12} \] \par Từ (11), suy ra \[ IA^2=OI^2-R^2. \] Kết hợp với (12), ta được \[ IA^2=ID^2. \] Suy ra \[ IA=ID. \] \par Vậy ta đã chứng minh được \[ OH\cdot OA=OT\cdot OI \quad \text {và} \quad IA=ID. \] \end {enumerate}  \hfill \qedEX 
\end{Solbook}
